\documentclass[../cuaderno6.tex]{subfiles}

\begin{document}
\subsection{Prisma vertical}
En el caso del prisma en su posición vertical, no se registran diferencias para el campo coercitivo $H_c$ ni para la magnetización máxima $M_{max}$ como se puede observar en las figuras \ref{fig:H_c_segun_H_a_Prisma_X_vertical_Fa0_24_09} y \ref{fig:M_max_segun_H_a_Prisma_X_vertical_Fa0_24_09} respectivamente, mientras que la magnetización remanente $M_r$, visible en la figura \ref{fig:M_r_segun_H_a_Prisma_X_vertical_Fa0_24_09} decae muy levemente al aumentar el ángulo.
\begin{figure}[H]
    \centering
    \includegraphics[width=0.5\textwidth]{../img/Fa0/Prisma X vertical/H_c_segun_H_a_Prisma X vertical_Fa0_24_09.png}
    \caption{El campo coercitivo $H_c$ en función del logaritmo del campo aplicado $H_{max}$ para cada ángulo.}
    \label{fig:H_c_segun_H_a_Prisma_X_vertical_Fa0_24_09}
\end{figure}
\begin{figure}[H]
    \centering
    \includegraphics[width=0.5\textwidth]{../img/Fa0/Prisma X vertical/M_r_segun_H_a_Prisma X vertical_Fa0_24_09.png}
    \caption{Magnetización remanente $M_r$ en función del campo aplicado $H_{max}$.}
    \label{fig:M_r_segun_H_a_Prisma_X_vertical_Fa0_24_09}
\end{figure}


\begin{figure}[H]
    \centering
    \includegraphics[width=0.5\textwidth]{../img/Fa0/Prisma X vertical/M_max_segun_H_a_Prisma X vertical_Fa0_24_09.png}
    \caption{Magnetización máxima $M_{max}$ en función del campo aplicado $H_{max}$.}
    \label{fig:M_max_segun_H_a_Prisma_X_vertical_Fa0_24_09}
\end{figure}

El resultado del ajuste para el campo coercitivo $H_c^0$, visible en la figura \ref{fig:H0_c_angulo_minors_Prisma_vertical} no presenta variaciones claras, y en la mayoría de los casos no presenta diferencias significativas. El factor de desmagnetización $N_d$ tiene un mínimo en $15^\circ$, pero luego de este sólo crece con el ángulo, como sería de esperar según los resultados de la magnetización remanente $M_r$.

\begin{figure}[H]
    \centering
    \includegraphics[width=0.4\textwidth]{../img/Fa0/Prisma X vertical/H0_c_angulo_minors_Prisma X vertical.png}
    \caption{Gráfico de la pendiente del ajuste $H_c^0$ en función del ángulo.}
    \label{fig:H0_c_angulo_minors_Prisma_vertical}
\end{figure}
\begin{figure}[h]
    \centering
    \includegraphics[width=0.4\textwidth]{../img/Fa0/Prisma X vertical/N_d_angulo_major_Prisma X vertical.png}
    \caption{Factor de desmagnetización $N_d$ en función del ángulo.}
    \label{fig:N_d_angulo_major_Prisma_X_vertical}
\end{figure}

\end{document}